\documentclass[11pt,letterpaper]{article}

\usepackage[margin=1in]{geometry}
\usepackage{amsmath,amssymb}
\usepackage{booktabs}
\usepackage{multirow}
\usepackage{hyperref}
\usepackage{microtype}
\usepackage{enumitem}

\hypersetup{
  colorlinks=true,
  linkcolor=black,
  citecolor=black,
  urlcolor=blue
}

\title{Adaptive Temporal Retrieval\\
for Time-Sensitive Question Answering}

\author{
Daksh Agarwal \quad Shashank Kumar\\
University School of Automation and Robotics (USAR)\\
Guru Gobind Singh Indraprastha University
}

\date{November 2026}

\begin{document}
\maketitle

\begin{abstract}
The pipeline can label a question as present, past, or neither, keep a date window around the asked time, and read a sentence whose date agrees with the question.
This paper reports what the saved runs change, including the full tests dated 18 November 2026, and where the measurements stop.

On the templated test of 77 questions, adaptive exact match is 73 of 77 and a fixed $\pm 2$ year filter (B3) is 35 of 77.
Containment is 76 of 77 against 75 of 77.
On 37 of the 38 questions that only the adaptive row matches, B3's prediction already contains the gold string and the adaptive string.
That exact-match gap is answer projection, which B2 and B3 do not run.
Precision at 5 is $0.650$ against $0.585$.
Skipping retrieval fires on 10 of those 77 questions, only when a classifier says ``no'' and the text hits a hardcoded list of trivia.
It does not make the system faster: adaptive mean latency is $1.491$ seconds and the median is $1.560$, against $1.380$ and $1.305$ for B3.

On 193 templated questions about six entity names, adaptive exact match is 147 of 193 against 49 of 193 for B3, and containment is 147 of 193 against 137 of 193.
Of the 98 exact-match gains, 88 are questions whose B3 string already contains the gold answer.
Retrieval is never skipped.
The method still gives the adaptive system an entity field the baselines do not, so this is not a pure comparison of date windows.
Mean latency is $0.905$ seconds against $0.709$; the medians are $0.779$ and $0.668$.

The 500-question TimeQA run was completed.
Exact match is 0 for B2, B3, and the adaptive system.
Containment is 148 of 500, 129 of 500, and 96 of 500.
Median latencies are $3.730$, $3.541$, and $3.585$ seconds.
The adaptive mean is $47.486$ seconds because one query took $20974.8$ seconds; the mean of the other 499 is $5.547$ seconds.
A 16-question slice is not this result.
There is still no ablation and no comparison with a published temporal retriever.
\end{abstract}

\section{Introduction}
\label{sec:intro}

A fact such as who holds an office is true on an interval.
Retrieving the nearest passage can return a true sentence about another year.
Retrieval-augmented generation grounds a generator in an external corpus \cite{lewis2020rag}, but the usual pipeline has no separate decision for ``now'' and ``in 2016''.

This project adds a temporal policy on top of one dense index:
a classifier with a two-way necessity head and a three-way time head;
a date window whose width depends on the cue;
a sufficiency loop that may widen the search;
and a reader that prefers a sentence whose date matches the question.
The project title asks two questions, when to retrieve and what to retrieve.
The second is what the code mostly implements.
The first is a whitelist.
Section~\ref{sec:when} says exactly when retrieval is skipped, and Section~\ref{sec:results} shows that this skip does not change latency and does not occur on the held-out or TimeQA runs, including the 18 November 2026 full tests.

We do not claim a new state of the art.
We do not claim that the adaptive window beats standard retrieval once the answer stage is equal, except where a table shows that directly.
The tables come from three different programs.
Checkpoint~A is an earlier full synthetic test and an earlier full held-out test.
Checkpoint~B adds interval parsing, a surname filter, and date-aware sentence selection, and it was run on the first 16 test questions of each dataset only.
The full tests dated 18 November 2026 run B1, B2, B3, and the adaptive system on all 77 templated questions, all 193 held-out questions, and all 500 TimeQA questions.
Numbers from one of these runs are not improvements over another.

The contribution is a description of that pipeline, the formulas it uses, and a reading of the saved numbers that does not pick the friendliest metric.
The 18 November full tests are in Section~\ref{sec:fulltest}.
Shared answer extraction, ablations, and a comparison to published temporal rankers are still missing.
Those gaps are listed in Section~\ref{sec:limits}.

\section{Related work}
\label{sec:related}

\paragraph{Retrieval-augmented generation.}
Lewis et al.\ \cite{lewis2020rag} condition a generator on passages from a dense index.
Dense passage retrieval showed that a dual encoder can replace a sparse first stage for open-domain questions \cite{karpukhin2020dpr}.
Atlas trains the retriever and the reader together \cite{izacard2023atlas}.
Self-RAG learns when a generation should retrieve and when it should critique the result \cite{asai2024selfrag}.
Adaptive-RAG chooses among no retrieval, one retrieval, and several retrievals according to question complexity \cite{jeong2024adaptiverag}.
That is the closest system in name.
It is not a temporal ranker, and we do not run it.
A reader should not treat the shared word ``adaptive'' as a comparison.

\paragraph{Time in retrieval, before RAG.}
Li and Croft put document age into a language model so that recency-seeking queries do not depend on a hand-tuned boost \cite{li2003time}.
Kanhabua and N{\o}rv{\aa}g compared mixture models that add a temporal similarity to text similarity, on a news archive with time-scoped judgments \cite{kanhabua2011ranking}.
Those papers already treat ``what time is relevant?'' as a ranking problem.
Our window is a heuristic descendant of that idea, not a replacement of those models, and we do not evaluate on their collections.

\paragraph{Time in question answering.}
TempQuestions evaluates time expressions against a knowledge graph \cite{jia2018tempquestions}.
TimeQA aligns time-evolving Wikidata facts with Wikipedia \cite{chen2021timeqa}.
The metric in that paper is exact match.
Fusion-in-Decoder fine-tuned on Natural Questions and TimeQA reaches 60.5 exact match on the easy test and 46.8 on the hard test; human exact match on the hard test is 87.0 \cite{chen2021timeqa}.
The abstract's ``about 46\%'' is the hard-test figure, not a single number for the whole benchmark.
SituatedQA shows that one question can have different answers at different times \cite{zhang2021situatedqa}.
StreamingQA and RealTime QA test adaptation as new text arrives \cite{liska2022streamingqa,kasai2023realtimeqa}.
Dhingra et al.\ condition a language model on a timestamp so that closed-book recall of a changing fact depends on the year \cite{dhingra2022talm}.
FreshQA asks models about facts that change quickly, including false premises \cite{vu2023freshqa}.
TempRALM re-ranks retrieved passages with a temporal score and masks passages dated after the query time \cite{gade2024tempralm}.
We cite these systems to locate the project.
We did not reimplement them.
Neither the 16-question slice nor the 500-question TimeQA run is comparable to any of their reported scores.
The TimeQA subset used here is chunks of the pages attached to the sampled questions, not their retrieval corpus.

\paragraph{Readers and the first-stage stack.}
SQuAD 2.0 requires a span and an abstention decision \cite{rajpurkar2018squad2}.
The learned path uses a frozen SQuAD 2.0 model and does not fine-tune it.
Embeddings are BGE-small \cite{xiao2024cpack}, searched with an exact FAISS inner-product index \cite{johnson2019faiss}.
A MiniLM cross-encoder trained on MS MARCO reorders the first-stage hits \cite{wang2020minilm,nguyen2016msmarco,reimers2019sbert}.
The router is DistilBERT \cite{sanh2019distilbert,devlin2019bert}.

\section{Problem setup}
\label{sec:setup}

The corpus is $\mathcal{D}=\{(d_i,\tau_i)\}$.
Here $\tau_i$ is the timestamp stored in the index.
It is not the only date in $d_i$.
On the templated corpus, $\tau_i$ is the date of the fact.
On the TimeQA subset, $\tau_i$ is the first date the tagger finds in the chunk, or 1 January 1970 when it finds none.
A question $q$ may also name a target interval.
The system returns a string $\hat{y}$ and a bit saying whether retrieval ran.

Answer strings are normalised as in SQuAD: lower case, punctuation removed, whitespace collapsed.
Exact match is $1$ when the strings are equal.
Token F1 is the harmonic mean of token precision and token recall, taking the best gold string when several are listed.
Containment is
\begin{equation}
\mathrm{Cont}(\hat{y}, y) =
\begin{cases}
1 & \text{if } y \subseteq \hat{y} \text{ or } \hat{y} \subseteq y,\\
0 & \text{otherwise.}
\end{cases}
\end{equation}
The second branch matters.
A short prediction inside a long gold string scores $1$, and so does a 400-character sentence that merely contains the gold span.
On TimeQA the heuristic answers are sentences, so containment mostly asks whether the selected sentence contains the gold text.
It is not an answer-quality metric.
We report exact match beside it for that reason.

If the question is marked unanswerable, or the gold string is empty, exact match, token F1, and containment are replaced by abstention accuracy:
$1$ when the system abstains on an unanswerable question, and $1$ when it does not abstain on an answerable one.
Abstention means an empty string or a short refusal such as ``insufficient evidence''.
A score of $0.250$ on 16 questions means the decision was right on 4 of them.

Precision and recall at $k{=}5$ use the gold document ids stored with the question.
Precision divides the number of gold ids in the retrieved list by the length of that list, capped at 5.
On the templated data the generator stores up to three passages of the named entity, preferring those nearest the asked date.
On TimeQA a gold id is a chunk that contains an answer string, and the list is capped at 8.
A precision of $0.673$ is therefore an average hit rate inside the top 5, not evidence that each question has three or four relevant documents.
Questions with an empty gold list are omitted from the average.

Temporal accuracy looks only at $\tau$ of the top document.
For a current question, that timestamp must fall within three years of 1 October 2025.
For a historical question, its year must fall within two of a year written in the question or the time hint.
The metric does not read $\hat{y}$.
It is undefined for TimeQA and is left blank there.
For the fixed-window baseline the metric restates the filter.
That baseline keeps documents whose stored year is within $\pm 2$ of the hinted year, and the metric asks whether the top stored year is within $\pm 2$.
If the filter keeps any document, every survivor passes, and so does the top document.
The score can fall below $1$ only when the filter is empty and the fallback document lies outside the band.
It is not independent evidence that the answer is about the right year.

\section{Method}
\label{sec:method}

The pipeline classifies the question, retrieves with the dense index, reorders with the cross-encoder, applies a temporal policy, checks sufficiency, and reads an answer.
The index and the cross-encoder are shared.
The reader is not fully shared, and neither is the subject filter.
Section~\ref{sec:exp} states the difference before any table.

\subsection{Query router}
\label{sec:router}

DistilBERT maps $q$ to the classification-token state $h$.
One linear head emits necessity logits $z^{\mathrm{nec}}\in\mathbb{R}^{2}$.
Another emits temporal logits $z^{\mathrm{tmp}}\in\mathbb{R}^{3}$, over current, historical, and timeless.
A third head emits a confidence logit $z^{\mathrm{conf}}$.
Training uses class weights $w$ and minimises
\begin{equation}
\mathcal{L}
= \mathrm{CE}\!\left(z^{\mathrm{nec}}, y^{\mathrm{nec}}; w^{\mathrm{nec}}\right)
+ \lambda\,
\mathrm{CE}\!\left(z^{\mathrm{tmp}}, y^{\mathrm{tmp}}; w^{\mathrm{tmp}}\right),
\label{eq:loss}
\end{equation}
with $\lambda = 1.75$.
The confidence head is not in $\mathcal{L}$.
Nothing trains its weights, so they stay at initialization.
The encoder still moves, because the two classification losses depend on it.

At inference the class logits are divided by a temperature $T$ before the softmax.
Let $p^{\mathrm{nec}}$ and $p^{\mathrm{tmp}}$ be those distributions, and let $c=\sigma(z^{\mathrm{conf}})$.
The scalar that the policy reads is
\begin{equation}
\gamma =
\begin{cases}
\tfrac12 c + \tfrac12 \max\{\max p^{\mathrm{nec}}, \max p^{\mathrm{tmp}}\}
& \text{when blending is on,}\\[0.3em]
\max\{\max p^{\mathrm{nec}}, \max p^{\mathrm{tmp}}\}
& \text{otherwise.}
\end{cases}
\label{eq:conf}
\end{equation}
The main analyzer blends.
Because $c$ was never trained, half of $\gamma$ is the output of an untrained sigmoid.
If $\gamma < 0.55$, the policy does not trust the temporal class and retrieves in semantic mode.
An untrained head can therefore change which ranker runs.
We did not count how often $\gamma$ falls below $0.55$ on the test sets.
The router used on the learned path turns blending off, so its $\gamma$ is the maximum class probability.

\subsection{When retrieval is skipped}
\label{sec:when}

Retrieval is skipped only when the necessity class is $0$ and the question string contains one of a fixed list of trivia cues: arithmetic such as ``2 + 2'', ``H2O'', ``days are in a week'', ``spell '', and a few similar items.
A necessity prediction of $0$ on any other question still retrieves, under the name of a soft timeless lookup.
Low confidence also retrieves.
The skip is a whitelist guarded by a classifier, not a learned abstention from retrieval.

There is no parametric model underneath the skip.
The closed-book control is the same extractive reader with an empty document list, plus a few hardcoded replies for those trivia items.
A skip cannot be scored as ``the model already knew the answer.''
It can only be scored as ``the index was not called.''

\subsection{First-stage retrieval}

Documents are embedded with BGE-small English v1.5.
Queries use that model's required prefix.
Search is FAISS \texttt{IndexFlatIP}, exact on these corpora.
Evaluation asks for 80 candidates.
Their FAISS scores are min--max normalised to $\tilde{s}^{\mathrm{f}}_i\in[0,1]$ inside the pool.
The cross-encoder score $s^{\mathrm{ce}}_i$ is left in its raw range.
The blend is
\begin{equation}
s_i = 0.35\,\tilde{s}^{\mathrm{f}}_i + 0.65\,s^{\mathrm{ce}}_i .
\label{eq:blend}
\end{equation}
The top 40 values of $s_i$ are kept.
The two terms do not share a scale, so $0.35$ and $0.65$ are not the weights of a mixture of comparable quantities.
A cross-encoder score of $6$ dominates a dense score in $[0,1]$ whatever the written weights say.
We did not normalise $s^{\mathrm{ce}}$ in the reported runs, and we did not choose the weights on a development split.
They are constants from implementation.

\subsection{Dates}
\label{sec:dates}

Dates are read by rules.
ISO dates and month-year expressions are kept.
A bare year is placed on 15 June, and that midpoint is dropped when the same year already has a fuller date.
Intervals cover $y$--$y$, ``from \ldots to \ldots'', ``until'' or ``through'', two-digit seasons such as 2004--05, and ``for $n$ years beginning in $y$'' with $n\le 40$.
A span longer than 45 years is discarded, which is aimed at birth--death pairs and will also discard a long reign or a long career.
``(born \ldots)'' in the question is stripped before the target year is chosen.
``Before month year'' becomes the interval from three years earlier through that month.
Relative phrases such as ``last year'', ``since 2010'', and ``the predecessor'' are not given their own target.
A decade cue is recognised only when the tagger emits one; there is no general ordinal logic.

Let $I(d)$ be the intervals in a passage and $U(d)$ the point dates, including a usable timestamp.
The distance from a target date $t$ to the passage is
\begin{equation}
\delta(d, t) =
\begin{cases}
0 & \text{if } t \in [u,v] \text{ for some } [u,v]\in I(d),\\[0.25em]
\min\limits_{x \in \partial I(d)\cup U(d)} |x-t| & \text{otherwise,}
\end{cases}
\label{eq:delta}
\end{equation}
where $\partial I(d)$ is the set of endpoints.
Distance zero means the asked date falls inside a stated range.

\subsection{Window and subject filter}

Historical questions use the half-widths in Table~\ref{tab:windows}.
While no surviving document satisfies $\delta(d,t)\le 365.25\,h$, the half-width grows by
\begin{equation}
h \leftarrow \min(h_{\max},\; 1.5h + 0.25).
\label{eq:grow}
\end{equation}
Growth stops at the first in-window survivor or at $h_{\max}$.
We did not ablate the early stop, so we cannot say that stopping early changes the answer.
The widths themselves were set by hand.

\begin{table}[t]
\centering
\caption{Initial and maximum half-width of the historical window, in years.}
\label{tab:windows}
\begin{tabular}{lcc}
\toprule
Cue & Initial $h$ & Maximum $h$ \\
\midrule
Exact calendar date & 0.35 & 1.5 \\
Early, mid, or late in a year & 1.5 & 3.5 \\
A bare year & 1.0 & 4.0 \\
A decade & 5.0 & 8.0 \\
Historical, but no year & 3.0 & 8.0 \\
\bottomrule
\end{tabular}
\end{table}

On the historical branch, before the window test, the pool is cut to a subject.
If any candidate has a non-empty entity field, only documents whose overlap with that field is within $0.01$ of the best overlap are kept.
The baselines do not make this cut.
On the templated corpus the field is the entity name, filled when the passage was generated, so part of the historical gain can come from that metadata rather than from the date window.
Current mode does not use this cut.
It keeps a lexical tie and sorts by timestamp, and the fixed-window baseline calls that same function when it finds no year.
The overlap text includes the entity string whenever the field is filled, so both systems see it there.
If the field is empty, as it is for every TimeQA chunk, a document is kept when lexical overlap is at least $0.05$ and the text contains a surname, unless overlap is within $0.01$ of the best document.
The surname is the last capitalised token of the first person-like span, ignoring months and question words, and it must be at least three characters.
The rule is English-specific.
It drops a pronoun or a given name, and it was added while inspecting failures on the same 16 TimeQA questions used in Table~\ref{tab:slice}.

After a stop list, if $H$ remaining query tokens occur in the document, the lexical boost is $b=\min(0.25, 0.05H)$.
The stop list includes months, years, and words such as ``between'', ``held'', and ``team'', because those tokens were matching the wrong passage.

\subsection{Temporal score}

Among documents that survive, semantic scores are min--max normalised to $\tilde{s}_i$.
For current questions, with reference date $r=\text{1 October 2025}$,
\begin{equation}
g^{\mathrm{cur}}_i = \max\left(0,\; 1 - \frac{\max(r-\tau_i,\, 0)}{3650}\right).
\label{eq:cur}
\end{equation}
The denominator is about ten years.
Current mode then keeps documents whose lexical overlap is within $0.01$ of the best and sorts them by $\tau_i$, latest first.
The displayed score of rank $i$ in that tie is $1-0.05\,i/n$, so an older tied document is not treated as weak evidence by the sufficiency check.
``Latest timestamp'' is a proxy for ``current holder'' only when the index date is the fact date.
The reference day is frozen, and it is more than a year before the date of this paper.
No live benchmark is used.

For historical questions,
\begin{equation}
g^{\mathrm{hist}}_i
= \max\left(0,\;
1 - \frac{\delta(d_i, t)}{\max(90,\; 365.25\,w)}
\right),
\label{eq:hist}
\end{equation}
with $w$ the window width in years.
On a tagged corpus, and only when the caller sets a time-first flag, rank is
\begin{equation}
f_i = g^{\mathrm{hist}}_i + 0.01\,\tilde{s}_i + b_i .
\label{eq:timefirst}
\end{equation}
The fixed-window baseline does not set that flag.
On untagged prose with the flag set, a shrunk semantic score is avoided by $\tilde{s}'_i=0.5+0.5\tilde{s}_i$ and
\begin{equation}
f_i = 0.45\,\tilde{s}'_i + 0.55\,g^{\mathrm{hist}}_i + b_i .
\label{eq:untagged}
\end{equation}
Otherwise
\begin{equation}
f_i = \alpha \tilde{s}_i + (1-\alpha)\,g_i\tilde{s}_i + b_i ,
\label{eq:gate}
\end{equation}
with $\alpha=0.6$.
The product $g_i\tilde{s}_i$ is the default, including for the fixed-window baseline's call into the same function.
The top $k=5$ documents go to the reader.
Tagged and untagged questions therefore do not use the same ranking formula.
Saying that one pipeline was applied to every dataset would hide that branch.

\subsection{Sufficiency}

During retrieval the judge sees the min--max semantic score of each document after temporal reordering, not the raw cross-encoder score.
Let $s_{(1)}$ be that score for the top-ranked document, $\overline{s}$ the mean, $C$ the fraction of content query tokens found in the retrieved text, $T\in[0,1]$ a temporal fit, and $b_{(1)}$ the lexical boost of the top document.
\begin{equation}
\begin{aligned}
S =\;& 0.35\,C + 0.25\,\min(1, s_{(1)}) + 0.15\,\overline{s}\\
&+ 0.15\,T + 0.10\,\min(1, b_{(1)}/0.25).
\end{aligned}
\label{eq:suff}
\end{equation}
Because the scores were min--max normalised inside the pool, $s_{(1)}=1$ whenever the temporally top document is also the semantic maximum, which is common.
The $0.25$ term is then a constant.
A pool of poor documents that still differ from each other receives the same saturated value.
$T$ takes the better of the stored timestamp and $\delta(d,t)$, so a passage whose index date is a birth year can still score well if the text contains the asked date.
Evidence is accepted during retrieval when the subject check passes and $S\ge 0.55$.
The weights and the threshold were not fit to data.

If the check fails, the loop may enlarge the candidate pool, raise $k$, add to the window, or fall back toward semantic ranking.
At most two repair rounds run.
After an answer is drafted, an abstention stub can trigger one semantic retry, and the retry replaces the answer only when its sufficiency score is strictly higher.
The reported tables do not record how often any of these repairs fired.
A higher latency on the held-out test is consistent with more work, and it is not a measurement of the loop.

\subsection{Reading the answer}
\label{sec:reader}

The offline reader does not call an external language model.
B1, B2, and B3 receive whatever this reader returns.
The adaptive system then runs an extra projection.

The shared reader has a few shortcuts.
Trivia cues have fixed replies.
Definition questions return the first sentence.
A passage that contains only one long sentence, and matches an office or winner template, is shortened to the name.
Initials are not sentence breaks, so ``B.\ R.\ Gavai'' stays one token.
That rule was added after a split initial hurt exact match on a templated question.

For a longer passage the reader scores sentences by date agreement, content-word overlap, and a bonus for higher-ranked documents.
A fragment shorter than 24 characters, including a section header, stays attached to the next sentence.
A dated clause shorter than 48 characters is prepended with the previous sentence.
The best sentence is truncated at 400 characters.
This selector is shared.
On the TimeQA slice it moves standard retrieval as well as the adaptive system.
It was developed against failures in that slice, including career headers and list sentences.

Projection, applied only after the adaptive reader, extracts a name from general who-question and result patterns such as ``held the office'' or ``$X$ won''.
B2 and B3 do not run it.
An exact-match gap that appears only for the adaptive row, while containment is tied, is this projection.
It is not a retrieval result.
Section~\ref{sec:fulltest} counts how often that pattern appears in the 18 November office tests.
The patterns do not reliably extract a title such as ``assistant district attorney'' from a sentence that mentions it in passing.

\subsection{Learned path}
\label{sec:learned}

The learned path is a separate stack.
It is not the system in the full-test tables.

Its router starts from the saved analyzer, freezes the bottom four of six DistilBERT layers, and trains for one epoch on 240 questions.
Temperature $T=2.8$ was chosen on 48 questions, and the accuracy we report is on those same 48 questions: necessity accuracy $0.812$ (F1 $0.889$) and temporal accuracy $0.812$ (macro-F1 $0.782$).
That is not a held-out estimate.
The drop from the original analyzer's validation accuracy of $1.0$ is also not a like-for-like change.
The original figure is on the templated validation split.
The $0.812$ figure is on this 48-question calibration split.
The original weight file was not overwritten.

The ranker is gradient boosting \cite{friedman2001gbm}, 80 trees of depth 3, on 357 rows of which 143 are positive.
Features include the cross-encoder and FAISS scores, lexical overlap, date distance, the router's class probabilities, and two products: current-probability times recency, and historical-probability times $1/(1+\text{year distance})$.
At inference a temporal prior
\begin{equation}
\pi_i
= p_{\mathrm{cur}}\,\rho_i
+ p_{\mathrm{hist}}\frac{1}{1+\Delta_i}
+ p_{\mathrm{time}}\,\tilde{s}^{\mathrm{ce}}_i
\label{eq:prior}
\end{equation}
is mixed as
\begin{equation}
f_i = \eta_i\big(0.35\,p^{\mathrm{gb}}_i + 0.60\,\pi_i + 0.05\,\tilde{s}^{\mathrm{ce}}_i\big),
\label{eq:learnedrank}
\end{equation}
where $\rho_i$ is a recency rank, $\Delta_i$ is year distance, and the subject gate $\eta_i$ is $1$ or $0.25$.
The prior, not the trees, carries most of the weight.

The reader is \texttt{deepset/tinyroberta-squad2}.
The span score is the product of the start and end softmax probabilities.
The no-answer score is $1-\sigma(m)$ with the margin $m$ clipped to $[-20,20]$.
A negative margin does not by itself reject the span.
On a multi-sentence passage the reader first sees the date-matched sentence, and a span scoring at least $0.02$ is kept; otherwise the top document is read if its span scores at least $0.05$.
A logistic regression on reader score, no-answer score, and cross-encoder score was fit on 16 labelled rows after standardising the features.
A hand override also keeps a span when the reader score is at least $0.08$ and the cross-encoder score is at least $1$.
This path retrieves 40 candidates and keeps 20, against 80 and 40 for the heuristic systems, so its latency is not comparable even before the model difference.

\section{Experiments}
\label{sec:exp}

\subsection{Systems, and what they do not share}

Five systems are named in the tables.
\begin{itemize}[leftmargin=1.4em]
  \item \textbf{B1.} The extractive reader with no documents. Hardcoded trivia replies still fire. This is not a closed-book language model.
  \item \textbf{B2.} Always retrieve. Rank by \eqref{eq:blend}. No date window and no entity cut. The shared sentence selector runs. Projection does not.
  \item \textbf{B3.} Always retrieve. If a year is found, keep candidates whose stored year is within $\pm 2$, and if that set is empty keep the unfiltered pool. Then call the shared ranker without the time-first flag. No entity cut. No projection.
  \item \textbf{Adaptive.} Router, whitelist skip, entity or surname cut, adaptive window, sufficiency loop, shared sentence selector, and projection.
  \item \textbf{Learned.} Section~\ref{sec:learned}. Reported only on 16-question slices.
\end{itemize}
B1 appears only in the full synthetic table, because that is the saved run that included it.

\subsection{Data}

The templated corpus is generated from timelines of office holders, chief executives, sports winners, awards, and a few laws.
Each fact is a short dated passage whose entity field is the timeline name.
The generator writes questions about the holder at a date, the current holder, and timeless trivia, then fills a target size with paraphrases that prepend a stock phrase to an existing question.
The main file has 803 passages and 3{,}000 questions: 2{,}849 train, 74 validation, and 77 test.
Gold document ids are up to three passages of that entity, nearest the asked date.
One test question asks who was president of Brazil in 2016.
The corpus has no passage for the holder in that year.
Every ranker that stays inside 2016 misses it.
That is a hole in the generated data.

The held-out file rebuilds the same generator with a target of 1{,}200 questions, then moves every question that names one of six entities into the test split.
The six names are every fifth entity in a sorted list: Academy Award for Best Picture, CEO of Apple, CEO of Twitter/X, GDPR, President of Brazil, and UEFA Champions League winners.
Test size becomes 193 because those questions are dumped into test, not because a 15\% split was drawn.
Train and validation of this file do not contain the six names.
The analyzer used at test time is the one trained on the main file, and the main file does contain them.
The split therefore does not hold the entities out of the router.
Almost every other component is a rule and was not trained on question labels.
Withholding names from a training split, while evaluating rules, does not test generalisation of a learned retriever.

The TimeQA subset is built from \texttt{hugosousa/TimeQA}.
The loader would take a half-easy, half-hard sample if rows carried a level field.
We did not record whether that field was present, so this paper does not claim an easy or a hard evaluation.
The built subset has 5{,}931 chunks and 900 questions, 400 train and 500 test.
Chunks are pieces of the page attached to the sampled question, not a retrieval corpus over Wikipedia.
A retrieved chunk is a much easier target than a passage from the open encyclopedia.
Gold ids are chunks containing the answer string, at most eight.

\subsection{Two checkpoints}

Checkpoint~A is the saved full synthetic run and the saved full held-out run reported in Tables~\ref{tab:syn77} and~\ref{tab:held193}.
It does not include the sentence selector, surname filter, or interval handling used in Checkpoint~B.
Checkpoint~B is the first 16 test questions of each dataset, in file order, after those changes.
The 16 questions are not a random sample.
The TimeQA rules in Checkpoint~B were edited against failures in that file order.
We do not treat Checkpoint~B as an unbiased estimate of the 500-question test.
The 18 November 2026 files are that full test, for all three datasets, and they are a different save from Checkpoint~A.
Adaptive containment on the TimeQA slice is 13 of 16; on the 500-question file it is 148 of 500.

Runs are on a laptop CPU.
No GPU was used.
An earlier draft of this paper said the 500-question TimeQA test had not been run, and it called that a hardware limit.
The 18 November files are the run.
One adaptive query took $20974.8$ seconds, so the wall-clock time was not the half hour that a four-second median would suggest.
The other three systems finished that same query in well under ten seconds.
Section~\ref{sec:fulltest} gives the mean, the median, and that query.

The reference date inside \eqref{eq:cur} and inside temporal accuracy is 1 October 2025, the date the templated exam was built against.
Current-question behaviour is measured on that templated exam only.

The main analyzer was trained on the templated training questions with six augmentations, 17{,}094 rows, batch size 8.
The saved log contains 7 epochs.
The configured budget was 15.
We do not have a record of why training stopped, and we did not check whether an augmentation of a validation question occurs in the training rows.
Validation necessity accuracy and temporal accuracy are both $1.0$.
On this data the label is often the word ``current'' or a written year.
The figure says the classifier fits those cues.
It does not say the classifier reads time in Wikipedia prose.
No real-text test of the router is reported.

\section{Results}
\label{sec:results}

\subsection{Checkpoint A, templated test}

Table~\ref{tab:syn77} is the full 77-question test from Checkpoint~A.
It is not the 18 November 2026 file in Table~\ref{tab:fulltest}.
Adaptive containment is $1.000$ and B3 containment is $0.987$: 77 questions against 76.
Exact match is $0.468$ for both, and token F1 is $0.676$ against $0.674$.
The abstract of an earlier draft led with containment.
Exact match, which is the stricter metric, does not move.
Precision at 5 does move, from $0.612$ to $0.673$, with recall at 5 unchanged at $0.954$.
That precision gap is the retrieval difference worth discussing, and it is confounded with the entity-field cut that only the adaptive system applies.

Retrieval rate $0.870$ is 67 of 77 questions.
The other 10 are the whitelist.
Mean latency is $0.56$ seconds for both Adaptive and B3.
Skipping those 10 questions does not show up as a faster system.
B2, which has no date logic, is worse on every answer metric.
That gap is real on this file and is the comparison a date filter is supposed to win.
It does not separate the adaptive window from the fixed two-year filter.

\begin{table}[t]
\centering
\caption{Checkpoint~A. Templated test, $n{=}77$. EM is exact match, F1 is token F1, Cont is containment, P@5 and R@5 use the gold document ids, Temp is the timestamp metric of Section~\ref{sec:setup}, Ret is the fraction of questions that called the index, and $s$ is mean seconds.}
\label{tab:syn77}
\small
\begin{tabular}{lcccccccc}
\toprule
System & EM & F1 & Cont & P@5 & R@5 & Temp & Ret & s \\
\midrule
B1 no retrieval & 0.117 & 0.158 & 0.130 & 0 & 0 & 0.247 & 0 & $\approx 0$ \\
B2 standard & 0.390 & 0.529 & 0.714 & 0.393 & 0.655 & 0.662 & 1 & 0.59 \\
B3 fixed $\pm 2$y & 0.468 & 0.674 & 0.987 & 0.612 & 0.954 & 1.000 & 1 & 0.56 \\
Adaptive & 0.468 & 0.676 & 1.000 & 0.673 & 0.954 & 1.000 & 0.870 & 0.56 \\
\bottomrule
\end{tabular}
\end{table}

Temporal accuracy is $1.000$ for both Adaptive and B3.
For B3 that is the circular reading in Section~\ref{sec:setup}.
No confidence interval is reported.
The pipeline is close to deterministic given the index, and the containment difference is one question.
A significance test would not turn one question into a stable finding.
We did not run one.

\subsection{Checkpoint A, six entity names}

Table~\ref{tab:held193} is the 193-question file from the same checkpoint, not Table~\ref{tab:fulltest}.
Adaptive exact match is $0.358$, against $0.285$ for B3 and $0.197$ for B2.
Containment is $0.834$, $0.741$, and $0.383$.
Token F1 follows: $0.587$, $0.500$, $0.276$.
Precision at 5 is $0.701$ against $0.614$ and $0.343$.
Recall at 5 ties B3 at $0.881$.
These gaps are larger than on the 77-question file.
They are still one checkpoint, without the sentence selector, with the entity field visible only to the adaptive ranker, and with a router that has seen the six names in the other training file.
Retrieval rate is $1$.
The ``when'' decision does nothing on this split.
Latency is about $1.02$ seconds against $0.80$ for B3.
The extra time is not a saving from skipping retrieval.

\begin{table}[t]
\centering
\caption{Checkpoint~A. Questions naming six entities, $n{=}193$. Columns match Table~\ref{tab:syn77}.}
\label{tab:held193}
\small
\begin{tabular}{lcccccccc}
\toprule
System & EM & F1 & Cont & P@5 & R@5 & Temp & Ret & s \\
\midrule
B2 standard & 0.197 & 0.276 & 0.383 & 0.343 & 0.572 & 0.596 & 1 & 0.79 \\
B3 fixed $\pm 2$y & 0.285 & 0.500 & 0.741 & 0.614 & 0.881 & 1.000 & 1 & 0.80 \\
Adaptive & 0.358 & 0.587 & 0.834 & 0.701 & 0.881 & 0.984 & 1 & 1.02 \\
\bottomrule
\end{tabular}
\end{table}

B3's temporal accuracy is again $1.000$, and Adaptive's is $0.984$.
The fixed filter matches its own metric slightly more often.
The answer metrics go the other way.
Both facts can be true, and neither is a reason to quote temporal accuracy as answer quality.

The containment--exact gap is large because many gold strings are short status phrases.
A sentence can name the right person and still miss the phrase.
Containment scores the phrase if it happens to occur.
That is a property of the metric and of the labels.

\subsection{Checkpoint B, first 16 questions}

Table~\ref{tab:slice} is Checkpoint~B.
It is a development slice.
We report it because it is the only measurement of the sentence selector, and because hiding it would make the TimeQA story look cleaner than it is.

On the templated and held-out slices, Adaptive and the learned path reach exact match $1.000$ and $0.938$.
B3 matches their containment and scores exact match $0$, because B3 returns the sentence and does not run projection.
Giving the baselines the same projection, or dropping exact match from this table, would remove that gap.
These 16 templated questions sit at the start of the file and are easier than the full test in Table~\ref{tab:syn77}, where exact match is $0.468$.
The learned path's $1.000$ and $0.938$ must not be placed next to the full-test numbers.

On TimeQA, before Checkpoint~B's date and sentence changes, containment on these same 16 questions was 4 of 16 for B2 and 1 of 16 for both B3 and Adaptive.
After the changes it is 12, 7, and 13 of 16.
B2 gains eight questions from a selector it shares.
Adaptive moves from 1 of 16 to 13 of 16.
It started three questions behind standard retrieval and finished one question ahead, after edits chosen on this slice.
That is not evidence that the window transferred to Wikipedia text.
Exact match stays $0$ for every heuristic row: the system returns a sentence, and the gold label is a span inside it.
Token F1 is $0.162$ for Adaptive.

The learned reader scores exact match $0.188$, three of 16, and containment $0.188$.
It is the worst row on this slice, below B3.
Its recall at 5 is the highest ($0.601$), so the retrieved set is not the whole failure.
Abstention accuracy is $0.250$, four of 16 correct decisions, because the 16-example logistic model abstains on questions whose span is present.
One of the 16 questions is unanswerable.
The heuristic system answers it from a plausible organisation, so its abstention accuracy is 15 of 16: it is right on the answerable items and wrong on the one item that should be refused.

B3 is weakest on this slice ($0.438$ containment) for a structural reason.
Its filter uses $\tau_i$, and on these chunks $\tau_i$ is often a birth year or the placeholder 1 January 1970.
The $\pm 2$ year cut removes the paragraph that contains the asked year.
This baseline is a poor control for untagged Wikipedia text.
On the templated corpus, where $\tau_i$ is the fact date, it is a much stronger one, which is why Table~\ref{tab:syn77} is nearly a tie.

\begin{table}[t]
\centering
\caption{Checkpoint~B. First 16 test questions, in file order, after the date and sentence changes. Not a test of the 500-question TimeQA split. Temp is omitted. Learned is Section~\ref{sec:learned}.}
\label{tab:slice}
\small
\begin{tabular}{llccccc}
\toprule
Data & System & EM & F1 & Cont & R@5 & s \\
\midrule
\multirow{4}{*}{TimeQA}
 & B2 & 0 & 0.148 & 0.750 & 0.499 & 3.19 \\
 & B3 & 0 & 0.078 & 0.438 & 0.190 & 3.11 \\
 & Adaptive & 0 & 0.162 & 0.812 & 0.573 & 3.98 \\
 & Learned & 0.188 & 0.188 & 0.188 & 0.601 & 4.84 \\
\midrule
\multirow{4}{*}{Held-out}
 & B2 & 0 & 0.154 & 0.312 & 0.458 & 0.62 \\
 & B3 & 0 & 0.438 & 0.938 & 1.000 & 0.51 \\
 & Adaptive & 0.938 & 0.938 & 0.938 & 0.979 & 0.58 \\
 & Learned & 0.938 & 0.938 & 0.938 & 1.000 & 0.41 \\
\midrule
\multirow{4}{*}{Synthetic}
 & B2 & 0 & 0.296 & 0.688 & 0.542 & 0.69 \\
 & B3 & 0 & 0.393 & 1.000 & 0.979 & 0.56 \\
 & Adaptive & 1.000 & 1.000 & 1.000 & 0.958 & 0.63 \\
 & Learned & 1.000 & 1.000 & 1.000 & 0.875 & 0.45 \\
\bottomrule
\end{tabular}
\end{table}

The learned path is faster on the office slices ($0.41$ and $0.45$ seconds against about $0.58$ and $0.63$).
It also retrieves a smaller pool.
The timing is not evidence that a trained ranker is cheaper at the same search depth.

\subsection{Full tests of 18 November 2026}
\label{sec:fulltest}

Table~\ref{tab:fulltest} is the later full-test save: all 77 templated questions, all 193 held-out questions, and all 500 TimeQA questions.
B1 is included because that save recorded it.
The learned path is not in these files.
Latency is the recorded per-query time.
The mean is the arithmetic mean.
The median is the midpoint of the sorted times, averaging the two central values when the count is even.
Answer metrics are rounded to three decimals in the table.
Exact-match and containment counts below are the unrounded totals.

\begin{table}[t]
\centering
\caption{Full tests dated 18 November 2026.
EM is exact match, F1 is token F1, Cont is containment, P@5 and R@5 use the gold document ids, and Temp is the timestamp metric of Section~\ref{sec:setup}.
Temp is undefined on TimeQA.
Ret is the fraction of questions that called the index.
Mean and Med.\ are seconds.
Scores other than the raw counts in the text are rounded to three decimals.
B1 mean and median latency are below $10^{-3}$ seconds on every row (synthetic means $4.25\times 10^{-5}$ seconds).}
\label{tab:fulltest}
\footnotesize
\setlength{\tabcolsep}{3.2pt}
\begin{tabular}{lccccccccc}
\toprule
System & EM & F1 & Cont & P@5 & R@5 & Temp & Ret & Mean & Med. \\
\midrule
\multicolumn{10}{l}{\emph{Templated, $n{=}77$}} \\
B1 no retrieval & 0.117 & 0.158 & 0.130 & 0 & 0 & 0.247 & 0 & 0.000 & 0.000 \\
B2 standard & 0.390 & 0.534 & 0.727 & 0.393 & 0.655 & 0.662 & 1 & 1.396 & 1.358 \\
B3 fixed $\pm 2$y & 0.455 & 0.662 & 0.974 & 0.585 & 0.954 & 1.000 & 1 & 1.380 & 1.305 \\
Adaptive & 0.948 & 0.960 & 0.987 & 0.650 & 0.954 & 1.000 & 0.870 & 1.491 & 1.560 \\
\midrule
\multicolumn{10}{l}{\emph{Six entities, $n{=}193$}} \\
B1 no retrieval & 0.000 & 0.002 & 0.000 & 0 & 0 & 0.000 & 0 & 0.000 & 0.000 \\
B2 standard & 0.197 & 0.287 & 0.409 & 0.343 & 0.572 & 0.596 & 1 & 0.719 & 0.664 \\
B3 fixed $\pm 2$y & 0.254 & 0.469 & 0.710 & 0.542 & 0.903 & 1.000 & 1 & 0.709 & 0.668 \\
Adaptive & 0.762 & 0.766 & 0.762 & 0.675 & 0.902 & 0.984 & 1 & 0.905 & 0.779 \\
\midrule
\multicolumn{10}{l}{\emph{TimeQA, $n{=}500$}} \\
B1 no retrieval & 0 & 0.014 & 0.014 & 0 & 0 & --- & 0 & 0.000 & 0.000 \\
B2 standard & 0 & 0.072 & 0.258 & 0.235 & 0.467 & --- & 1 & 3.726 & 3.541 \\
B3 fixed $\pm 2$y & 0 & 0.061 & 0.192 & 0.154 & 0.281 & --- & 1 & 3.746 & 3.585 \\
Adaptive & 0 & 0.079 & 0.296 & 0.241 & 0.488 & --- & 1 & 47.486 & 3.730 \\
\bottomrule
\end{tabular}
\end{table}

On the templated file, adaptive exact match is 73 of 77 and B3 is 35 of 77.
Containment is 76 of 77 against 75 of 77, and recall at 5 is tied at $0.954$.
The exact-match gap is not a retrieval gap.
B2 and B3 do not run the projection in Section~\ref{sec:reader}; the adaptive row does.
On 38 questions adaptive matches and B3 does not.
On 37 of those 38, B3's prediction already contains the gold string and already contains the adaptive string, and B3's containment bit is 1.
Containment is tied on 76 of 77 questions.
One illustration is the question ``As of now, who is the Prime Minister of India'', whose gold string is ``Narendra Modi''.
B3 returns ``Narendra Modi (Prime Minister of India)'' and scores exact match 0.
The adaptive row returns ``Narendra Modi'' and scores exact match 1.
Median prediction length on this file is 13 characters for the adaptive system and 50 for B3.
The same comparison against B2 accounts for 25 of 43 adaptive-only exact matches; the other B2 misses are questions B2's sentence does not contain.
Precision at 5 still moves, from $0.585$ to $0.650$.
That gap is confounded with the entity-field cut that Section~\ref{sec:method} gives only to the adaptive system.
The saved rows do not record the cut itself.
Retrieval rate $0.870$ is 67 of 77.
The other 10 are the whitelist.
Mean latency is $1.491$ seconds against $1.380$, and the medians are $1.560$ against $1.305$.
Skipping those 10 questions does not show up as a faster system.

On the six-entity file the same reading holds, with a small remainder that is not projection.
Adaptive exact match and containment are both 147 of 193.
B3 exact match is 49 of 193 and B3 containment is 137 of 193.
Of the 98 questions that adaptive alone matches, 88 already have the gold string, and the adaptive string, inside the B3 prediction.
Containment differs on 10 of 193 questions.
Median prediction length is 14 characters for the adaptive system and 23 for B3.
Retrieval rate is 1.
The ``when'' decision does nothing on this split.
Mean latency is $0.905$ seconds against $0.709$ for B3; the medians are $0.779$ and $0.668$.
The extra time is not a saving from skipping retrieval.
B3's temporal accuracy is again $1.000$ (193 of 193) and the adaptive score is $0.984$ (190 of 193).

On TimeQA, exact match is 0 of 500 for B1, B2, B3, and the adaptive system.
The predictions stay sentences: the median length is 144 characters for B2 and 140 for both B3 and the adaptive row, so the office-set projection is not what these strings show.
Containment is 148 of 500 for the adaptive system, 129 of 500 for B2, and 96 of 500 for B3.
B3 remains the weak control on this corpus, for the timestamp reason in the slice discussion.
Precision at 5 is $0.241$, $0.235$, and $0.154$.
Recall at 5 is $0.488$, $0.467$, and $0.281$.
Abstention accuracy is $0.820$ for all four systems, so it does not separate them.
The saved router necessity and temporal accuracies are $0.987$ (76 of 77) and $1.000$ on the templated test, and $1.000$ for both on the held-out and TimeQA tests.
This draft does not break those scores down by cue.
A written year is often present in a TimeQA question, and the templated labels are often the word ``current'' or a year, so these figures still do not show that the router reads time in prose that lacks a cue.

Latency on TimeQA has to be read as a median and a mean together.
B2's mean is $3.726$ seconds and its median is $3.541$.
B3's mean is $3.746$ and its median is $3.585$.
The adaptive median is $3.730$ seconds, close to those two.
The adaptive mean is $47.486$ seconds.
One query accounts for the gap.
It is \texttt{tqq\_c61f3b840e35}, ``Janet Frame received which award in 1990?'', and the recorded time is $20974.779$ seconds.
The saved policy label is \texttt{historical\_adaptive\_window}.
The returned sentence contains the gold string ``Order of New Zealand'' and does not exact-match it.
B3 contains that gold string as well, in $4.737$ seconds; B2 does not, in $3.523$ seconds.
The file does not record whether a repair loop ran.
The mean of the other 499 adaptive questions is $5.547$ seconds, still above B2 and B3.
Fifty adaptive queries took longer than 10 seconds, and four took longer than 30 seconds ($20974.8$, $39.775$, $32.616$, and $30.368$).
No B2 or B3 query on this file exceeded 10 seconds (maxima $8.676$ and $9.573$).
The 16-question slice, at about four seconds and at 13 of 16 containment, is not this distribution.

\subsection{Where the remaining TimeQA errors come from}

On the 16-question slice the heuristic misses are of three kinds.
Some articles never state the asked month, so the date scorer has nothing to prefer.
Some name the right year inside a list of teams, and a neighbouring sentence wins.
The unanswerable item is answered anyway.
None of these counts is a population rate.
They describe 16 questions that the rules have already seen.

\section{What these measurements do not show}
\label{sec:limits}

The sentences below are the limits of this draft.
They are not future work written as a compliment to the method.

Checkpoint~A, the 16-question slice, and the 18 November 2026 full tests are different programs.
Table~\ref{tab:fulltest} is the only measurement here of all 77, all 193, and all 500 questions together.
It is not a claim that the code at the end of Section~\ref{sec:method} was frozen and re-run under one configuration file.
The slice is not an estimate of the 500-question test: adaptive containment there is 13 of 16, and on the full file it is 148 of 500.
Checkpoint~A ties adaptive and B3 exact match at $0.468$ on 77 questions.
The later templated full test does not reproduce that tie.

No component was removed in these files.
The router, the early stop in \eqref{eq:grow}, the sufficiency loop, the surname cut, and the sentence selector each have a story and no ablation.
On the slice, the TimeQA movement that survives after subtracting B2 is one question, and the shared selector accounts for the rest of B2's jump from 4 of 16 to 12 of 16.
On the 500-question file the adaptive row contains the gold string on 148 questions and B2 on 129, a gap of 19 questions, with no ablation that isolates the window.

The baselines do not receive projection.
The adaptive system does.
On the office full tests that is most of the exact-match gap: 37 of 38 templated questions and 88 of 98 held-out questions where adaptive matches and B3 does not are questions whose B3 string already contains the answer.
Exact match on those sets is not a retrieval comparison.
The baselines also do not receive the entity-field cut.
Precision on the templated sets is not a pure comparison of date policies.
Containment on the templated full test is almost a tie with B3 (76 of 77 against 75 of 77).
On the held-out full test it differs by 10 questions (147 against 137).

B3 is a reasonable control when the stored timestamp is the fact date.
It is a weak control on TimeQA, and its perfect temporal accuracy is largely the filter scoring itself.
There is no row for B2 plus a date-aware re-ranker, B2 plus the frozen SQuAD reader, Adaptive-RAG, or TempRALM.

The skip decision was not tested against a model that knows the trivia without the whitelist.
Latency does not improve.
Held-out and TimeQA never skip.
The project title's first half is not a result.

Constants throughout Section~\ref{sec:method} --- the blend weights, $\lambda$, $\alpha$, the window table, the 45-year cutoff, the 15 June midpoint, the $0.55$ threshold, the $0.01$ tie --- were chosen by hand.
Some of the TimeQA rules were chosen while looking at the test slice.
There is no untouched development split for those choices.

The learned path is a frozen reader, a tree model on 357 rows, and a logistic model on 16 rows plus an override.
On the only Wikipedia slice it is the worst system.
Its router accuracy was computed on the split used to pick $T$.

Containment is bidirectional and, for this reader, close to ``did the gold string occur in the chosen sentence?''
Temporal accuracy does not read the answer, and the same words are used for the router's class accuracy in Section~\ref{sec:learned}.
The two quantities are not the same.

The current-question branch is scored against a frozen day on templated text.
FreshQA and RealTime QA were not run.
A deployment clock and this exam clock are different.

We do not release a public code URL with this draft.
The implementation is the project tree that produced the tables.

\section{Conclusion}

The pipeline can prefer a dated sentence and can keep a subject-matched window.
On the 18 November 2026 templated test, containment is one question ahead of a fixed two-year filter, 76 of 77 against 75 of 77.
Exact match is 73 of 77 against 35 of 77, and 37 of those 38 extra matches are questions where the fixed-filter string already contained the gold answer.
That is answer projection, which B2 and B3 do not run, not a retrieval result.
The held-out file shows the same pattern on 88 of 98 exact-match gains, with containment ahead by 10 questions (147 of 193 against 137).
Precision at 5 is higher on both office files, under a comparison that also gives the adaptive system an entity field.
Skipping retrieval is a trivia whitelist, used on 10 of 77 questions, unused on the other two full tests, and invisible in the latency.
Means and medians on the office files are higher for the adaptive system than for B3.

On the 500-question TimeQA run, exact match is 0 for every heuristic system.
Containment is 148 of 500 against 129 of 500 for standard retrieval and 96 of 500 for the fixed filter.
The adaptive median is $3.730$ seconds.
The adaptive mean is $47.486$ seconds because one query, ``Janet Frame received which award in 1990?'', took $20974.8$ seconds and still did not beat B3 on that item.
The 16-question slice, where the adaptive row finishes one question ahead of standard retrieval after rules written on those same questions, is not this result.
The learned reader is not in the full-test files.
A paper that wants to claim a temporal policy on real pages still has to give every baseline the same projection and the same metadata, and to ablate the window.
This draft does not contain those experiments.

\section*{Acknowledgments}
Dr.\ Annu Priya supervised this project at the University School of Automation and Robotics, Guru Gobind Singh Indraprastha University.

\begin{thebibliography}{99}

\bibitem{lewis2020rag}
P.~Lewis, E.~Perez, A.~Piktus, F.~Petroni, V.~Karpukhin, N.~Goyal, H.~K\"uttler, M.~Lewis, W.~Yih, T.~Rockt\"aschel, S.~Riedel, and D.~Kiela.
Retrieval-augmented generation for knowledge-intensive NLP tasks.
In \emph{Advances in Neural Information Processing Systems}, 2020.

\bibitem{karpukhin2020dpr}
V.~Karpukhin, B.~Oguz, S.~Min, P.~Lewis, L.~Wu, S.~Edunov, D.~Chen, and W.~Yih.
Dense passage retrieval for open-domain question answering.
In \emph{Proceedings of EMNLP}, 2020.

\bibitem{izacard2023atlas}
G.~Izacard, P.~Lewis, M.~Lomeli, L.~Hosseini, F.~Petroni, T.~Schick, J.~Dwivedi-Yu, A.~Joulin, S.~Riedel, and E.~Grave.
Atlas: Few-shot learning with retrieval augmented language models.
\emph{Journal of Machine Learning Research}, 24(251):1--43, 2023.

\bibitem{asai2024selfrag}
A.~Asai, Z.~Wu, Y.~Wang, A.~Sil, and H.~Hajishirzi.
Self-RAG: Learning to retrieve, generate, and critique through self-reflection.
In \emph{International Conference on Learning Representations}, 2024.

\bibitem{jeong2024adaptiverag}
S.~Jeong, J.~Baek, S.~Cho, S.~J.~Hwang, and J.~C.~Park.
Adaptive-RAG: Learning to adapt retrieval-augmented large language models through question complexity.
In \emph{Proceedings of NAACL}, 2024.

\bibitem{li2003time}
X.~Li and W.~B.~Croft.
Time-based language models.
In \emph{Proceedings of CIKM}, pages 469--475, 2003.

\bibitem{kanhabua2011ranking}
N.~Kanhabua and K.~N{\o}rv{\aa}g.
A comparison of time-aware ranking methods.
In \emph{Proceedings of SIGIR}, 2011.

\bibitem{jia2018tempquestions}
Z.~Jia, A.~Abujabal, R.~Saha~Roy, J.~Str\"otgen, and G.~Weikum.
TempQuestions: A benchmark for temporal question answering.
In \emph{Companion of The Web Conference}, 2018.

\bibitem{chen2021timeqa}
W.~Chen, X.~Wang, and W.~Y.~Wang.
A dataset for answering time-sensitive questions.
In \emph{NeurIPS Datasets and Benchmarks}, 2021.
arXiv:2108.06314.

\bibitem{zhang2021situatedqa}
M.~J.~Q.~Zhang and E.~Choi.
SituatedQA: Incorporating extra-linguistic contexts into QA.
In \emph{Proceedings of EMNLP}, 2021.

\bibitem{liska2022streamingqa}
A.~Li\v{s}ka, T.~Ko\v{c}isk\'y, E.~Gribovskaya, T.~Terzi, E.~Sezener, D.~Agrawal, C.~de~Masson~d'Autume, T.~Scholtes, M.~Zaheer, S.~Young, E.~Gilsenan-McMahon, S.~Austin, P.~Blunsom, and A.~Lazaridou.
StreamingQA: A benchmark for adaptation to new knowledge over time in question answering models.
In \emph{Proceedings of ICML}, 2022.

\bibitem{kasai2023realtimeqa}
J.~Kasai, K.~Sakaguchi, Y.~Takahashi, R.~Le~Bras, A.~Asai, X.~Yu, D.~Radev, N.~A.~Smith, Y.~Choi, and K.~Inui.
RealTime QA: What's the answer right now?
In \emph{NeurIPS Datasets and Benchmarks}, 2023.

\bibitem{dhingra2022talm}
B.~Dhingra, J.~R.~Cole, J.~M.~Eisenschlos, D.~Gillick, J.~Eisenstein, and W.~W.~Cohen.
Time-aware language models as temporal knowledge bases.
\emph{Transactions of the Association for Computational Linguistics}, 10:257--273, 2022.

\bibitem{vu2023freshqa}
T.~Vu, M.~Iyyer, X.~Wang, N.~Constant, Jerry~Wei, Jason~Wei, C.~Tar, Y.-H.~Sung, D.~Zhou, Q.~Le, and T.~Luong.
FreshLLMs: Refreshing large language models with search engine augmentation.
arXiv:2310.03214, 2023.

\bibitem{gade2024tempralm}
A.~Gade and J.~Jetcheva.
It's about time: Incorporating temporality in retrieval augmented language models.
arXiv:2401.13222, 2024.

\bibitem{rajpurkar2018squad2}
P.~Rajpurkar, R.~Jia, and P.~Liang.
Know what you don't know: Unanswerable questions for SQuAD.
In \emph{Proceedings of ACL}, 2018.

\bibitem{xiao2024cpack}
S.~Xiao, Z.~Liu, P.~Zhang, N.~Muennighoff, D.~Lian, and J.-Y.~Nie.
C-Pack: Packaged resources to advance general Chinese embedding.
In \emph{Proceedings of SIGIR}, 2024.
arXiv:2309.07597.

\bibitem{johnson2019faiss}
J.~Johnson, M.~Douze, and H.~J\'egou.
Billion-scale similarity search with GPUs.
\emph{IEEE Transactions on Big Data}, 7(3):535--547, 2019.

\bibitem{wang2020minilm}
W.~Wang, F.~Wei, L.~Dong, H.~Bao, N.~Yang, and M.~Zhou.
MiniLM: Deep self-attention distillation for task-agnostic compression of pre-trained transformers.
In \emph{International Conference on Learning Representations}, 2020.

\bibitem{nguyen2016msmarco}
T.~Nguyen, M.~Rosenberg, X.~Song, J.~Gao, S.~Tiwary, R.~Majumder, and L.~Deng.
MS MARCO: A human generated machine reading comprehension dataset.
In \emph{CoCo@NeurIPS}, 2016.

\bibitem{reimers2019sbert}
N.~Reimers and I.~Gurevych.
Sentence-BERT: Sentence embeddings using Siamese BERT-networks.
In \emph{Proceedings of EMNLP}, 2019.

\bibitem{sanh2019distilbert}
V.~Sanh, L.~Debut, J.~Chaumond, and T.~Wolf.
DistilBERT, a distilled version of BERT: Smaller, faster, cheaper and lighter.
arXiv:1910.01108, 2019.

\bibitem{devlin2019bert}
J.~Devlin, M.-W.~Chang, K.~Lee, and K.~Toutanova.
BERT: Pre-training of deep bidirectional transformers for language understanding.
In \emph{Proceedings of NAACL}, 2019.

\bibitem{friedman2001gbm}
J.~H.~Friedman.
Greedy function approximation: A gradient boosting machine.
\emph{Annals of Statistics}, 29(5):1189--1232, 2001.

\end{thebibliography}

\end{document}
